\documentclass[11pt,final,hidelinks]{article}
\usepackage{lmodern}
\usepackage[T1]{fontenc}
\usepackage[english]{babel}
\usepackage[utf8]{inputenc}
\usepackage[activate={true,nocompatibility},final,tracking=true,kerning=true,
    spacing=true,factor=1100,stretch=10,shrink=10]{microtype}
\microtypecontext{spacing=nonfrench}
\usepackage[margin=1in]{geometry}
\usepackage{graphicx}
\graphicspath{{img/}}
\usepackage{caption}
\usepackage{subcaption}
\usepackage{booktabs}
\usepackage{array}
\usepackage{mathtools}
\usepackage{enumitem}
\usepackage{commath}
\usepackage{appendix}
\renewcommand\appendixtocname{Appendices}
\usepackage{amsmath}
\usepackage{amsthm}
\usepackage{amssymb}
\usepackage{pgfplots}
\usepackage{tikz}
\usepackage{tikzscale}
\usepackage[square,numbers]{natbib}
\bibliographystyle{plain}
\usepackage{siunitx}
\numberwithin{equation}{section}
\usepackage[hypertexnames=false]{hyperref}
\usepackage{cleveref}

\pgfplotsset{compat=newest}

\usepackage[fancy,section]{alex}

\hypersetup{
    pdftitle={Composition of Bernoulli sets},
    pdfauthor={Alexander Towell},
    pdfsubject={computer science},
    pdfkeywords={
        probabilistic data structure,
        approximate set,
        bloom filter,
        set operations,
        interval arithmetic,
        compositional algebra},
    colorlinks=true,
    linkcolor=magenta,
    citecolor=green,
    filecolor=blue,
    urlcolor=green
}

\title
{
    Composition of Bernoulli sets
}
\author
{
    Alexander Towell\\
    \texttt{atowell@siue.edu}
}
\date{2026}

\begin{document}
\maketitle
\begin{abstract}
Bernoulli sets---random approximate sets with independent, element-wise errors parameterized by false positive and false negative rates---are closed under set-theoretic composition.
We derive closed-form error rates for complement, union, intersection, and set difference, showing that every such operation on Bernoulli sets yields a (possibly higher-order) Bernoulli set whose error rates are computable functions of the operand rates and set cardinalities.
The resulting structure is monoidal under union and intersection, supports exact computation of relational predicates (subset and equality probabilities), and admits an interval arithmetic extension that propagates uncertain rate parameters through set operations with guaranteed bounds.
These results build on the Bernoulli set model introduced in a companion paper~\cite{bernoulliSets}, which establishes the axioms, error-count distributions, and asymptotic limits assumed here.
\end{abstract}

\microtypesetup{protrusion=false}
\tableofcontents
\microtypesetup{protrusion=true}

%%%%%%%%%%%%%%%%%%%%%%%%%%%%%%%%
% Introduction
%%%%%%%%%%%%%%%%%%%%%%%%%%%%%%%%

\section{Introduction}

The Bernoulli set model~\cite{bernoulliSets} provides a probabilistic framework for \emph{random approximate sets}---sets that approximate an objective set with independent, element-wise errors parameterized by false positive and false negative rates.
That companion paper establishes the axioms, derives the binomial distributions of error counts, and develops the asymptotic normal limits and confidence intervals for the error rates.

This paper addresses the \emph{compositional} question: what happens when Bernoulli sets are combined through set-theoretic operations?
Practical systems rarely use a single approximate set in isolation.
Encrypted search indexes, database query planners, and distributed systems routinely compose approximate sets through union, intersection, complement, and difference.
Without a compositional algebra, each new combination requires a bespoke error analysis.

\paragraph{Contribution.}
We show that Bernoulli sets are \emph{closed} under set-theoretic composition: every operation on Bernoulli sets yields a (possibly higher-order) Bernoulli set whose error rates are closed-form functions of the operand rates and set cardinalities.
Specifically, we derive:
\begin{enumerate}[nosep]
    \item closed-form error rates for complement, union, intersection, and set difference, together with the monoidal structure they induce,
    \item exact formulas for the relational predicates (subset and equality probabilities) on Bernoulli sets, and
    \item an interval arithmetic extension that propagates uncertain rates through set operations, yielding guaranteed bounds on output error rates.
\end{enumerate}

\paragraph{Organization.}
\Cref{sec:preliminaries} restates the definitions and results from~\cite{bernoulliSets} needed in this paper.
\Cref{sec:set_theory} derives error rates for all four set-theoretic operations and establishes closure properties.
\Cref{sec:relational} treats relational predicates (subset, equality) on Bernoulli sets.
\Cref{sec:intervals} develops interval arithmetic for uncertain rate parameters.

%%%%%%%%%%%%%%%%%%%%%%%%%%%%%%%%%
% Preliminaries
%%%%%%%%%%%%%%%%%%%%%%%%%%%%%%%%%

\section{Preliminaries}
\label{sec:preliminaries}

We briefly restate the elements of the Bernoulli set model~\cite{bernoulliSets} needed in this paper.

\subsection{Bernoulli set model}

Let $U$ be a finite universal set with $u = |U|$, and let $A \subseteq U$ be an objective set.
A \emph{Bernoulli set} $\tilde{A}$ is a random approximate set of $A$ governed by two axioms.

\begin{axiom}[Element-wise independence {\cite[Axiom~1]{bernoulliSets}}]
\label{asm:element_indep}
For all distinct $x, y \in U$, the error events at $x$ and $y$ are mutually independent.
Furthermore, within each partition block $B_i$ (\cref{def:model_order}), the error indicators $\{\mathbf{1}_{\tilde{A}}(x) : x \in B_i\}$ are identically distributed.
\end{axiom}

\begin{axiom}[Conditional independence of block error rates {\cite[Axiom~2]{bernoulliSets}}]
\label{asm:fpr_fnr_r_indep}
Given the objective set $R$, the random error rates $\alpha_1, \ldots, \alpha_n$ for the respective partition blocks $B_1, \ldots, B_n$ are mutually independent.
\end{axiom}

The Bernoulli set $\tilde{A}$ is parameterized by its expected false positive rate $\fprate$ and expected true positive rate $\tprate$, writing $\tilde{X}[\fprate][\tprate]$.
The complementary rates are $\fnrate = 1 - \tprate$ (FNR) and $\tnrate = 1 - \fprate$ (TNR).

\begin{proposition}[Element-level error probabilities {\cite[Prop.~1]{bernoulliSets}}]
\label{prop:element_rates}
By \cref{asm:element_indep}, for any single element $x \notin A$,
$\Prob{\mathbf{1}_{\tilde{A}}(x) = 1} = \fprate$, and for $x \in A$,
$\Prob{\mathbf{1}_{\tilde{A}}(x) = 1} = \tprate$.
\end{proposition}

\begin{definition}[Model order {\cite[Def.~3]{bernoulliSets}}]
\label{def:model_order}
The \emph{order} of a Bernoulli set model is the number of partition blocks $n$.
A zeroth-order model ($n = 0$) is exact.
A first-order model ($n = 1$) assigns a single error rate (symmetric channel).
A second-order model ($n = 2$) distinguishes positives from negatives with rates $\tprate$ and $\fprate$.
In general, an $n$-th order model partitions $U$ into $n$ blocks, each with its own error rate.
\end{definition}

\subsection{Distributions and asymptotic limits}

The number of false positives given $n$ negatives is binomially distributed: $\FP_n \sim \bindist(n, \fprate)$.
Similarly, the number of false negatives given $p$ positives satisfies $\FN_p \sim \bindist(p, \fnrate)$.

The random false positive rate $\alpha_n = \FP_n / n$ converges in distribution to
\begin{equation}
\label{eq:fpr_clt}
    \alpha_n \xrightarrow{d} \normdist(\fprate, \fprate(1-\fprate) / n),
\end{equation}
and the random true positive rate $\TPR_p$ converges to $\normdist(\tprate, \tprate(1-\tprate) / p)$.
Asymptotic $\gamma \cdot 100\%$ confidence intervals for the false positive rate and true positive rate are respectively
\begin{equation}
\label{eq:conf_fpr}
\fprate \pm \sqrt{\frac{\fprate(1-\fprate)}{n}} \Phi^{-1}\!\left(\frac{1+\gamma}{2}\right)
\end{equation}
and
\begin{equation}
\label{eq:conf_tpr}
\tprate \pm \sqrt{\frac{\tprate(1-\tprate)}{p}} \Phi^{-1}\!\left(\frac{1+\gamma}{2}\right),
\end{equation}
where $\Phi^{-1}$ is the inverse CDF of the standard normal.
See~\cite{bernoulliSets} for proofs.

\input{sections/set_theory}
\input{sections/relational}
\input{sections/uncertain_rates}

%%%%%%%%%%%%%%%%%%%%%%%%%%%%%%%%%
% Conclusion
%%%%%%%%%%%%%%%%%%%%%%%%%%%%%%%%%

\section{Conclusion}
\label{sec:conclusion}

We have shown that Bernoulli sets are closed under all four set-theoretic operations: complement, union, intersection, and set difference.
Each operation yields a (possibly higher-order) Bernoulli set whose error rates are closed-form functions of the operand rates and set cardinalities.
The resulting collection of Bernoulli sets forms commutative monoids under union and intersection, with the empty set and universal set serving as identity elements.
Relational predicates---subset and equality---admit exact probability formulas that factorize element-wise.
An interval arithmetic extension propagates uncertain rate parameters through set operations, yielding guaranteed bounds when parameters are only approximately known.

\paragraph{Practical implications.}
A practitioner building a system on Bloom filters need not re-derive the error analysis for each new combination of filters: the composition theorems of \cref{sec:set_theory} apply immediately, yielding closed-form error rates and confidence intervals for any set-theoretic expression.

\paragraph{Open problems.}
The partition weights $w_1, w_2, w_3$ in the set-operation formulas require knowledge of the objective set cardinalities, which are typically unknown (\cref{rem:unknown_partitions}).
Developing principled estimators for these cardinalities from the approximate sets themselves, with quantified bias and variance, is a natural next step.

\paragraph{Companion papers.}
The Bernoulli set model, including the axioms, error-count distributions, asymptotic limits, confidence intervals, and higher-order composition theorem, is developed in~\cite{bernoulliSets}.
The probability distributions of binary classification measures (positive predictive value, negative predictive value, accuracy, Youden's $J$) are derived in~\cite{bernoulliMeasures}.
The joint entropy of error counts and the information-theoretic space complexity lower bound are developed in~\cite{bernoulliEntropy}.


\bibliography{references}

\end{document}
